\chapter{Introduction}
\label{chap:introduction}

\section{Background and Motivation}
\label{sec:intro:motivation}

Steel trusses remain one of the most material-efficient structural
forms for spanning long distances, supporting heavy roof and floor
loads, and carrying transmission and communication infrastructure.
Their design is dominated by member sizing: once topology and
connectivity are fixed by architectural and serviceability constraints,
the engineer is left with the problem of choosing cross-sectional areas
that minimise structural weight while satisfying stress, displacement,
slenderness, and stability limits. Because steel is roughly half of
the embodied cost and a comparable share of the embodied carbon of
typical steel-framed construction, even single-digit reductions in
weight translate directly into measurable economic and environmental
savings over the lifetime of a structure. Sizing optimisation is
therefore a perennially active research area in structural engineering
and a natural test bed for new algorithmic ideas.

The canonical formulation of the sizing problem is a constrained
nonlinear programme in which the objective is structural weight, the
design variables are member areas drawn from continuous or discrete
catalogues, and the constraints are evaluated by a finite-element (FEM)
solve. Schmit and Farshi~\cite{schmit1974} formalised the problem in
its modern form; subsequent benchmark studies by Sunar and
Belegundu~\cite{sunar1991}, Schmit and Miura~\cite{schmit1976}, and
Erbatur \emph{et al.}~\cite{erbatur2000} established the four small to
mid-scale benchmarks --- 10-bar planar, 25-bar spatial, 72-bar tower,
and 200-bar planar --- that today serve as the de facto comparison set
for any new optimisation method applied to truss sizing. These
problems are reproduced verbatim in the present work, and their
published optima are used as the validation reference throughout
\cref{chap:results}.

Classical evolutionary algorithms~--- genetic algorithms (GA), particle
swarm optimisation (PSO)~\cite{kennedy1995pso}, and the
non-dominated sorting genetic algorithm~II (NSGA-II)
\cite{deb2002nsga2} --- are the dominant solvers in the truss-sizing
literature. They are robust, gradient-free, and handle the
combinatorial structure of discrete-section catalogues without
modification. They have one well-documented weakness: each candidate
design requires a full FEM solve, and convergence on even a moderately
sized benchmark routinely consumes tens of thousands of evaluations.
For the four benchmarks studied here, a single 500-generation,
100-individual run takes between several minutes (10-bar) and several
hours (200-bar) on a workstation-class CPU. When ten random seeds are
required to estimate variance, and three algorithms are compared, the
total wall-clock cost becomes the practical bottleneck of any thorough
study.

Three families of modern artificial-intelligence techniques offer
complementary remedies to this bottleneck, and their integration into a
single coherent framework is the central contribution of this thesis.
First, neural surrogate models trained on FEM-generated data can
replace the inner-loop solve with a forward pass through a small
multi-layer perceptron, reducing per-evaluation cost by roughly two
orders of magnitude~\cite{queipo2005,forrester2008}. Second,
reinforcement-learning (RL) agents trained against the surrogate
environment can amortise design knowledge into a policy network that
proposes near-feasible designs in a single inference call, providing an
alternative to population-based search~\cite{schulman2017ppo,
hayashi2020rl,brown2022rl}. Third, large language models (LLMs) such
as Anthropic's Claude can interpret natural-language design briefs and
return an initial area vector that warm-starts the evolutionary loop,
bringing semantic engineering knowledge into a formerly purely
numerical pipeline~\cite{wei2022cot,geng2024,makatura2024}.

The Indian engineering context adds a final dimension. Domestic steel
construction is governed by IS~800:2007~\cite{is800-2007}, whose
slenderness, tension, compression, and serviceability clauses differ in
both form and conservatism from the AISC and Eurocode 3 provisions
typically referenced in international truss-optimisation literature
\cite{subramanian2008,duggal2014}. A framework that targets Indian
practice must therefore embed IS 800 checks directly into the
optimisation constraint set rather than treating them as a
post-processing audit. Doing so changes which designs are admissible
and, in our experience, frequently changes which design the optimiser
reports as best.

\section{Problem Statement}
\label{sec:intro:problem}

Given a planar or spatial pin-jointed steel truss with fixed topology,
material, and applied static load case, determine the vector of
member cross-sectional areas that minimises total structural weight
subject to stress, displacement, slenderness, and IS~800:2007 design
constraints, using an artificial-intelligence-accelerated evolutionary
framework whose every component is independently validated against
published benchmark data and whose end-to-end results are reproducible
from seeded random states.

\section{Research Objectives}
\label{sec:intro:objectives}

The thesis pursues five concrete objectives, each tied to a specific
validation gate developed in \cref{chap:methodology} and reported
against in \cref{chap:results}.

\begin{enumerate}
  \item \textbf{Build and validate a custom FEM engine} for planar and
    spatial pin-jointed trusses, verified to machine precision against
    the analytical 3-bar problem and used as the ground-truth solver
    for every downstream layer.

  \item \textbf{Implement classical optimisation baselines}
    (GA, PSO, NSGA-II) within a unified problem interface and
    demonstrate convergence to within published tolerance
    ($\pm 2$--$3\%$) on all four benchmarks across multiple random
    seeds.

  \item \textbf{Train a neural surrogate model} of the FEM solver that
    achieves a coefficient of determination $R^{2} > 0.98$ on weight
    prediction for held-out designs, and that yields at least a
    fifty-fold reduction in inner-loop evaluation time when substituted
    into the GA loop.

  \item \textbf{Develop and evaluate alternative AI design routes:}
    a proximal-policy-optimisation (PPO) reinforcement-learning agent
    trained on the surrogate environment, and an LLM warm-start
    interface that converts natural-language design briefs into initial
    area vectors. Each route is evaluated against the classical
    baseline using statistically meaningful protocols
    (paired \emph{t}-tests, $p < 0.05$).

  \item \textbf{Embed IS~800:2007 compliance} directly in the
    optimisation constraint set --- slenderness (Cl. 3.8), tension
    yielding and rupture (Cl. 6.2--6.3), compression (Cl. 7.1), and
    deflection (Cl. 5.6.1) --- and verify that every reported optimum
    passes a clause-by-clause audit and a hand-calculation spot check.
\end{enumerate}

A sixth, transverse objective is reproducibility: every numerical
result reported in \cref{chap:results} is generated by a documented
script, seeded explicitly, and reproducible end-to-end from a
Streamlit + FastAPI demo shipped with the thesis source.

\section{Scope and Limitations}
\label{sec:intro:scope}

The framework deliberately restricts its scope along several axes in
order to maintain depth where it matters.

\paragraph{Analysis type.} Only linear static analysis is considered.
Geometric and material nonlinearity, dynamic loading, seismic response,
and fatigue are out of scope. The static assumption is consistent with
the four benchmark formulations adopted here and is the standard
setting in the published truss-sizing literature.

\paragraph{Design problem.} Only sizing (cross-sectional area)
optimisation is performed. Topology and shape optimisation, while
active research areas in their own right, are explicitly excluded; the
truss connectivity is treated as fixed input.

\paragraph{Member behaviour.} All members are assumed pin-jointed and
carry axial force only. Connection design (bolts, welds), gusset
plates, and lateral-torsional buckling of bending members are out of
scope. Buckling is treated through IS~800 compression strength
provisions on individual members.

\paragraph{Design code.} The code-compliance layer implements
IS~800:2007 only. AISC 360 and Eurocode 3 implementations are
straightforward extensions but are not pursued here.

\paragraph{Benchmarks.} Validation is restricted to four canonical
problems --- 10-bar, 25-bar, 72-bar, and 200-bar --- chosen because
their geometries, loads, and reference optima are unambiguously
published and have been reproduced by dozens of independent groups.
Real industrial structures are not analysed.

\paragraph{Compute envelope.} All training and optimisation are
designed to complete on a standard laptop or workstation CPU. No GPU
acceleration is assumed, and no cloud compute is required to reproduce
any reported number.

\section{Organisation of the Thesis}
\label{sec:intro:organisation}

\Cref{chap:literature} surveys the relevant literature in three
strands --- classical evolutionary truss optimisation, neural
surrogate modelling for structural analysis, and emerging applications
of reinforcement learning and large language models to engineering
design --- and identifies the integrative gap that this thesis aims to
fill.

\Cref{chap:methodology} develops the eight-layer architecture in
detail: the FEM engine and its validation, the four benchmark
encodings, the GA / PSO / NSGA-II wrappers built on
pymoo~\cite{blank2020pymoo} with penalty-based constraint handling
\cite{coello2002}, the surrogate dataset generation and network
architecture, the PPO environment and reward shaping, the LLM prompt
and warm-start protocol, and the IS~800 compliance checker. Each
section concludes with the validation gate that subsequent chapters
report against.

\Cref{chap:results} presents the numerical results across all four
benchmarks and all algorithms. The 10-bar planar truss is treated
in greatest depth because it is the most heavily reported benchmark in
the literature and serves as the primary validation case. Surrogate
accuracy and speedup, PPO training behaviour, the LLM warm-start study,
the multi-objective Pareto fronts produced by NSGA-II, and the
IS~800:2007 clause-by-clause compliance verification are reported in
turn.

\Cref{chap:conclusion} summarises the contributions, places the
quantitative results in context against the published literature,
discusses the limitations exposed by the experimental work, and
identifies directions for future research --- including extension to
topology optimisation, integration with real Indian steel-section
catalogues per SP~6(1)~\cite{sp6-1-1964}, and migration of the
surrogate and LLM components to in-house or domain-specialised models.
